%versi 2 (8-10-2016) 
\chapter{Pendahuluan}
\label{chap:intro}
   
\section{Latar Belakang}
\label{sec:label}

\textit{Website} SharIF-Judge merupakan \textit{website} utama bagi mahasiswa informatika Universitas Katolik Parahyangan untuk berlatih dan mengasah kemampuan pemrograman. \textit{Website} ini memuat berbagai \textit{assignment} dan untuk masing-masing \textit{assignment} memiliki \textit{problem} yang harus diselesaikan mahasiswa sebagai bahan latihan. Saat ini \textit{website} SharIF-Judge melakukan penilaian secara otomatis dengan mencocokan keluaran kode dengan jawaban yang sudah disiapkan oleh pembuat soal dan memberikan nilai berdasarkan keakuratan keluaran yang dihasilkan oleh kode mahasiswa. Selain nilai yang diberikan oleh \textit{webiste} SharIF-Judge, dosen juga harus menilai kode secara manual untuk menilai logika dan efisiensi kode dari setiap mahasiswa. Namun evaluasi secara manual ini memiliki kendala tersendiri, dosen harus mengunduh dan memeriksa berkas jawaban mahasiswa satu per satu secara manual.
\dtext{5-10}

\section{Rumusan Masalah}
\label{sec:rumusan}
\begin{enumerate}
     \item Bagaimana mengimplementasikan pembatasan akses berdasarkan \textit{IPv4} dan penanda \textit{quiz} dalam satu \textit{instance}?
     \item Bagaimana merancang perubahan tampilan antarmuka agar memiliki status visual berupa emoji dan durasi pengerjaan setiap \textit{assignment} diberikan warna sesuai dengan status \textit{assignment}?
     \item Bagaimana cara melakukan pengunduhan berkas \textit{assignment} yang sudah dikompresi dalam bentuk zip dan memiliki penamaan \texttt{\{id-assignment\}-\{nama-assignment\}-\{nama-soal\}.zip}.
\end{enumerate}
\dtext{6}

\section{Tujuan}
\label{sec:tujuan}
\begin{enumerate}
     \item Mengimplementasikan pembatasan akses berdasarkan \textit{IPv4} dan penanda \textit{quiz} dalam satu \textit{instance}.
     \item Mengimplementasikan perubahan tampilan antarmuka agar memiliki status visual berupa emoji dan durasi pengerjaan setiap \textit{assignment} diberikan warna sesuai dengan status \textit{assignment}.
     \item Mengimplementasikan pengunduhan berkas \textit{assignment} yang sudah dikompresi menjadi bentuk zip dan memiliki penamaan \texttt{\{id-assignment\}-\{nama-assignment\}-\{nama-soal\}.zip}.
\end{enumerate}
\dtext{7}

\section{Batasan Masalah}
\label{sec:batasan}
Untuk mempermudah pembuatan template ini, tentu ada hal-hal yang harus dibatasi, misalnya saja bahwa template ini bukan berupa style \LaTeX{} pada umumnya (dengan alasannya karena belum mampu jika diminta membuat seperti itu)

\dtext{8}

\section{Metodologi}
\label{sec:metlit}
\begin{enumerate}
    \item Melakukan studi mengenai pembatasan akses \textit{website} berdasarkan IPv4.
    \item Melakukan studi mengenai cara mengunduh banyak file secara bersamaan dan dijadikan format zip.
    \item Merancang sistem untuk pembatasan akses berdasarkan IPv4 dan mengunduh banyak file lalu mengubahnya dalam bentuk zip.
    \item Mengimplementasikan perubahan sistem pada \textit{website} SharIF-Judge.
    \item Melakukan Pengujian dan eksperimen.
    \item Menganalisis hasil eksperimen.
    \item Menulis dokumen tugas akhir.
\end{enumerate}
\dtext{9}

\section{Sistematika Pembahasan}
\label{sec:sispem}
\begin{itemize}
    \item \textbf{Bab 1:} Pendahuluan \\
          Membahas latar belakang, rumusan masalah, tujuan, batasan masalah, metodologi, dan sistematika pembahasan.
    \item \textbf{Bab 2:} Landasan Teori \\
          Membahas teori-teori yang berhubungan dengan penelitian ini, yaitu SharIF-Judge, CodeIgniter 3, Twig, IDE, subnetting \textit{IPv4}, mengunduh dan melakukan format file.
    \item \textbf{Bab 3:} Analisis \\
          Membahas analisis terhadap perangkat lunak SharIF-Judge dan IDE pada SharIF-Judge.
    \item \textbf{Bab 4:} Perancangan \\
          Membahas perancangan fitur yang akan diimplementasikan pada SharIF Judge.
    \item \textbf{Bab 5:} Implementasi dan Pengujian \\
          Membahas implementasi fitur pada SharIF Judge dan pengujian yang dilakukan beserta dengan analisis hasil pengujian.
    \item \textbf{Bab 6:} Kesimpulan dan Saran \\
          Membahas kesimpulan dari penelitian ini dan saran untuk penelitian berikutnya.
\end{itemize}
\dtext{10}